\section{Lean formalization: definitions, verified core, and completion path}
\label{sec:formalization}

\subsection{The endpoint is an embedded proof}

The included standalone kernel reproduces three current repository files:
\texttt{Syntax.lean}, \texttt{Deduction.lean}, and \texttt{PR.lean}.
The source bodies were checked against the pinned commit. Original imports
are replaced by \texttt{import Std} and a natural-number notation; the
declarations and proofs are preserved. The archive retains the original
source files, their provenance and the repository license.

The key interfaces, in namespace \texttt{LeanProofs.ConciseIPRA}, are:
\begin{lstlisting}
PR k                 -- a fixed primitive recursive function symbol
PTm n                -- a term with n free variables
Fm PTm n             -- a first-order arithmetic formula
Prf PRAAx Gamma phi  -- an object-iPRA derivation
PRA.qfInd            -- quantifier-free object induction
\end{lstlisting}
Here \texttt{PR} is an inductive syntax with constructors for zero,
successor, projection, composition and primitive recursion. A constructor
of type \texttt{PR k} supplies a fixed program symbol, not an interpreter
for arbitrary programs. \texttt{Prf} is the intuitionistic natural
deduction relation with equality and the stated axiom scheme.

The current repository proves the exact transport statement
\begin{lstlisting}
pra_iff_native :
  (Gamma |- [PRAAx] phi) <->
  NPrf (Gamma.map NFm.ofPRA) (NFm.ofPRA phi)
\end{lstlisting}
where the displayed ASCII turnstile is only explanatory notation for
the repository's scoped judgment \cite{repoBridge}. The forward direction
is applied to an already constructed arithmetic derivation. It does not
accept a host theorem about naturals as a replacement. The standalone
build does not rebuild the full native equivalence dependency closure.

\subsection{What is checked in the supplied Lean core}

\begin{longtable}{P{.27\linewidth}P{.43\linewidth}P{.21\linewidth}}
\toprule
Module & Contents & Logical status\\
\midrule\endhead
{\footnotesize\path{ObjectKernel.lean}} & Pinned scoped syntax, substitution,
intuitionistic deduction and ordinary iPRA axioms. & Reused foundations.\\[3pt]
{\footnotesize\path{RankedDescent.lean}} & Run validity, rank budget, termination in the
initial rank, invariant transport and a contraction-to-descent lemma.
& Proved host theorems.\\[3pt]
{\footnotesize\path{CoherentTower.lean}} & Iteration of specified local lifts,
admissibility, adjacent coherence, base projection and uniqueness over
the seed. & Proved host theorems.\\[3pt]
{\footnotesize\path{ObjectComposition.lean}} & Witness-contract composition,
specialization, existential introduction, a QF-induction interface and
an instantiated identity-iteration proof. & Actual object-proof
constructions.\\[3pt]
{\footnotesize\path{DescentTargets.lean}} & Explicit PR codes for arithmetic and bounded
iteration; exact object formulas and the remaining descent compiler
goal. & Definitions; the compiler goal is unproved.\\[3pt]
{\footnotesize\path{ApplicationTargets.lean}} & Polynomial systems, Jacobians, ring and
filtration laws, frozen iteration, jets, residues and projections.
& Compiled definitions; general lifting and presentation goals unproved.\\[3pt]
{\footnotesize\path{WeightedNormTargets.lean}} & Rational polynomials, monic reduction,
multiplication matrices, determinants, weights and monomial counts.
& Compiled definitions; general norm and counting goals unproved.\\[3pt]
{\footnotesize\path{HurwitzTargets.lean}} & Binomial products, typed derivative demand,
and a scalar rational polynomial-ODE prefix producer.
& Compiled definitions; general algebra and ODE goals unproved.\\[3pt]
{\footnotesize\path{ApplicationChecks.lean}} & A modular root and a finite basis count;
exact rational jet and norm computations. & Two finite kernel theorems;
rational diagnostics are executable checks.\\[3pt]
{\footnotesize\path{AxiomAudit.lean}} & Printed dependencies of the core declarations.
& Reproducible audit.\\
\bottomrule
\end{longtable}

The core build uses Lean 4.32.0 and \texttt{Std}. Its coherence and
uniqueness declarations use no axioms. The other audited core results
use at most Lean's metatheoretic \texttt{propext} and \texttt{Quot.sound};
none uses \texttt{Classical.choice} or \texttt{sorryAx}. Those dependencies
are properties of the host proofs, not extra axioms inserted into
\texttt{PRAAx}. The build and axiom logs list the exact declarations.

All ten modules compile. The application propositions are ordinary
\texttt{Prop} definitions, not axioms or admitted proofs. Their compilation
checks that the definitions and targets are well typed; it does not prove
those targets. The modular example proves by kernel reduction that the
producer for \(X^3-X-2\), with \(q=4\), \(N=3\), seed \(2\) and inverse
\(3\), returns \(22\). The rational diagnostics check the Catalan prefix
\((1,1,2,5,14)\), a mixed-curve norm \(1+X+X^4\), and the Hurwitz
solution prefix \((1,1,2,6,24,120,720,5040)\) for \(Y'=Y^2\).
The concrete ODE target uses a scalar polynomial with constant rational
coefficients. The broader time-dependent, multivariate human theorem
extends the same finite coefficient recurrence; that extension is part
of the documented completion route.

\subsection{A genuine composition theorem}

For formulas \(P(x)\) and \(Q(y,x)\), the actual structure
\texttt{TermContract P Q} contains a fixed term \(f(x)\) and an object
derivation of \(Q(f(x),x)\) from \(P(x)\). A second contract contains
\(g(y,x)\) and a derivation of \(R(g(y,x),y,x)\) from \(Q(y,x)\).
The checked composition operation substitutes \(f(x)\) for \(y\), then
uses cut to produce
\[
        [P(x)]\ \vdash_{\ipra}\ R(g(f(x),x),f(x),x).
\]
Capture avoidance and context substitution are handled by the repository's
scoped syntax. The record contains proof obligations; there is no axiom
asserting that arbitrary preconditions and postconditions can be filled.

The file also defines the explicit PR symbol
\begin{lstlisting}
identityIteration : PR 2 :=
  .prec (.proj (0 : Fin 1)) (.proj (1 : Fin 3))
\end{lstlisting}
and proves in the embedded theory that iterating the identity on a seed
returns the seed. Its base uses \texttt{PRA.precZero} and
\texttt{PRA.proj}; its step uses \texttt{PRA.precSucc}, the accumulator
projection and the induction hypothesis. The outer induction is the
actual \texttt{PRA.qfInd}. This small example checks the full syntactic
route without misrepresenting the scale of the remaining algebra library.

\subsection{Numerical definitions for the application layers}

The following definitions are sufficient; no abstract analytic or
algebraic-geometric object is required as an object-theory primitive.

\begin{description}
\item[Ring and filtration.] A fixed finite tuple of PR symbols provides
\(\Valid(x)\), \(\operatorname{eq}(x,y)\), \(0,1,+,-,\cdot\) and
\(\operatorname{member}(N,k,x)\). The ring and filtration laws are
object derivations or explicit hypotheses. Their parameters include the
chosen precision and, for residues, the base \(q\).
\item[Polynomial system.] A polynomial is a list of pairs
\((c,e)\), where \(e\) is a vector of nonnegative exponents. A bounded
test verifies the vector length and coefficient validity. Evaluation
multiplies powers along each vector and then adds terms; formal partial
derivatives replace \(cX^e\) by
\(ce_jX^{e-\mathbf e_j}\) when \(e_j>0\). The Jacobian is the finite
matrix of evaluations of these partial derivatives.
\item[Jets.] A precision-\(N\) element is a coefficient vector of length
\(N\), with zero-default lookup. Equality compares the entries. Products
are truncated convolution; membership in level \(k\) checks that the
first \(\min(k,N)\) entries are zero. Hurwitz jets insert the natural
factor \(\binom nj\) in each product summand.
\item[Residues.] A precision-\(N\) element is the canonical natural
number below \(q^N\). Operations reduce signed integer expressions
modulo that modulus. The membership test is divisibility by
\(q^{\min(k,N)}\). Matrices and vectors are finite lists of these codes.
\item[Plane polynomials.] The supplied target module uses a finite list
of rational \(X\)-coefficient lists indexed by \(Y\)-degree, with
zero-default lookup. Equivalently, an arithmetic encoding can use sparse
triples \((i,j,a_{ij})\), normalized by adding duplicates and dropping
zeros. The monic remainder is computed by eliminating the largest
\(Y\) exponent. Conversion between the two finite representations is PR.
\item[Multiplication norm.] Construct column \(j<m\) by reducing
\(GY^j\). The matrix is an \(m\)-square array of rational
\(X\)-polynomials. Its determinant is either the signed permutation sum
or a recursively expanded Laplace determinant. Both are finite PR
algorithms; their equivalence can be proved coefficientwise.
\item[Point lists and subspaces.] A point is a pair of rational codes.
The common-zero predicate is evaluation equality. Pairwise distinctness
is the bounded conjunction over pairs of list indices. A subspace is a
finite rational spanning matrix; membership is a rank comparison, and
properness is rank below the ambient dimension. The weight-bounded
monomial list is enumerated by \(j<m\), \(i\le L\) and \(mi+dj\le L\).
\end{description}

Malformed arrays are explicitly excluded by bounded well-formedness tests.
The totalized operations may return defaults on malformed inputs. The
algebraic laws and result statements are always guarded by those tests.

\subsection{From the human lifting proof to a complete Lean proof}

Let \(H(N,P,U,a)\) be the bounded input test consisting of validity,
the residual membership and the two inverse-Jacobian congruences in
Theorem~\ref{thm:filtered-lifting}. Let \(B(N,P,U,a)\) be the state
after \(N-1\) updates. Its correctness target is
\[
 H(N,P,U,a)\to
 \bigl(P(B)=0\ \land\ B-a\in F^1\bigr).
\]
The uniqueness target additionally takes a \emph{fixed numerical candidate}
\(z\) and concludes
\[
 H(N,P,U,a)\land P(z)=0\land z-a\in F^1\quad\to\quad z=B.
\]
Every part is a bounded formula once arrays are encoded. Universal closure
over the input codes can be introduced after its proof.

The finite dependency order is as follows.
\begin{enumerate}
\item Prove the integer, rational and array operations correct on valid
codes. Give their \texttt{PR} syntax and their defining equations.
\item Prove addition and matrix multiplication preserve the indicated
filtration levels. The invariant is a conjunction over the processed
entries and the accumulated sum.
\item Prove polynomial evaluation, partial differentiation and the finite
Taylor identities. Induct on monomial exponents and on the length of
the processed term list; the invariant is a finite coefficient identity
and membership test, not a universally quantified ring identity.
\item Derive the single-step residual improvement equation (1) of the
lifting proof. Freeze \(N\), Booleanize invariant (3), and apply
\texttt{PRA.qfInd} to the explicit PR state recursion.
\item Freeze the two candidate roots and \(N\). Begin the uniqueness
induction at filtration zero, use the seed congruence for the step to
level one, and equation (2) for all later steps. At level \(N\), use
the proved nilpotence test.
\item Prove finite evaluation commutes with restriction. Compose this
with uniqueness to obtain pairwise coherence. This last step needs no
infinite family object.
\item Instantiate the generic contracts for ordinary jets, Hurwitz jets
and residues. For the Hurwitz ODE use a separate prefix recurrence;
prove coefficient causality by induction on the finite polynomial
evaluation, then induct on the stored solution length.
\end{enumerate}

At the end, replace each bounded formula by its fixed Boolean compiler
output and build terms of
\texttt{Prf PRAAx [] (allClose liftingFormula)}. The actual compiler
and the large algebraic derivation are not completed in this archive.
The definitions, local identities and induction formulas above make the
remaining work explicit rather than introducing an uninterpreted
``lifting soundness'' axiom.

\subsection{From the norm proof to a complete Lean proof}

Define \(\operatorname{GoodCurve}(m,d,c,F)\) to check
\(m\ge2\), \(d\ge1\), \(\gcd(m,d)=1\), \(c\ne0\), the two leading
coefficients and every lower-weight restriction. Define
\(\operatorname{ReducedNonzero}(m,G)\) by its finite coefficient tests,
and \(W(m,d,G)\) by bounded maximum over the nonzero terms. The main
object formula is the universal closure of
\[
 \begin{split}
 &\operatorname{GoodCurve}(m,d,c,F)\land
   \operatorname{ReducedNonzero}(m,G)\land
   \operatorname{DistinctCommonZeros}(F,G,L)\\
 &\hspace{35mm}\longrightarrow\quad
   \operatorname{length}(L)\le W(m,d,G).
 \end{split}
\]
Unlike a formula quantifying over a completed function field, this is an
ordinary finite assertion about the supplied polynomial and point codes.

\begin{longtable}{P{.28\linewidth}P{.34\linewidth}P{.29\linewidth}}
\toprule
Lemma to formalize & Finite mechanism & Induction invariant\\
\midrule\endhead
Monic remainder & Eliminate the highest \(Y\)-coefficient using \(F\).
& Original polynomial equals quotient times \(F\) plus residual;
processed high coefficients are zero.\\[3pt]
Norm specialization & Evaluate every entry in the multiplication matrix;
evaluate its finite determinant. & Equality of the accumulated
permutation sums.\\[3pt]
Rescaling identity & Clear a computed common power of \(t\), then compare
finite polynomial coefficients. & The change-of-basis equality for
processed matrix columns.\\[3pt]
Exact norm degree & Compute the cyclic companion determinant and compare
the constant term. & Finite coefficient support and a certified nonzero
top coefficient.\\[3pt]
Fiber multiplicity & Use the explicit Vandermonde rows and factor
\(X-u\) out of their entries. & The row identities and accumulated
factorization.\\[3pt]
Product of fiber factors & Remove distinct-coordinate powers successively.
& \(N_G=C_kQ_k\), with \(C_k\) the processed product, and the finite
Taylor coefficients needed for the next fiber.\\[3pt]
Dimension count & Enumerate the bounded monomials and residues modulo \(m\).
& Prefix counts and the permutation of residue classes.\\[3pt]
Subspace fiber bound & Row-reduce the rational spanning matrix and produce
an annihilator. & Row equivalence, nonzero pivots and the explicit free
variable giving a nonzero kernel vector.\\
\bottomrule
\end{longtable}

For all coefficient identities, the finite support bound is part of the
state. Monic reduction lowers the largest remaining \(Y\)-exponent; in
the principal weighted family it also never increases monomial weight.
The determinant has \(m!\) permutation terms. Rational values may grow
rapidly, but the state is still a PR-coded finite list. There is no need
for a common external bound on every possible rational input value.

An ambient algebra proof can be completed first using ordinary Lean
polynomials and matrices; that would be useful independent validation.
For the intended iPRA theorem, it must then be implemented in the numerical
coefficient library or accompanied by an explicit verified compiler that
constructs the corresponding embedded derivation. The target modules in
this archive make no general reflection assumption.

\subsection{Repository integration and reproducible builds}

The standalone source is intended for inspection and independent builds.
Inside ProveIt, replace the import of \texttt{ObjectKernel} by the current
\texttt{ConciseIPRA.PR} import and build against the pinned toolchain.
Do not import both versions of the same namespace in one build. Once a
new object theorem is completed, apply \texttt{pra\_iff\_native.mp} to
obtain its native form. The provided scripts compile the standalone
sources, run the axiom audit and regenerate exact-arithmetic checks.
The proof-status ledger distinguishes tested programs, compiled target
definitions, host proofs and embedded proof constructions.
